\input figfm1

\section{Background/Introduction}

This project will develop a {\bf foundation-model-driven, agentic-AI framework}
\cite{Vinuesa2026XAI} for the prediction and optimization of thermal-fluid
systems in fission and fusion reactors.
  Phase I demonstrator problems, illustrated in Fig.~\ref{fig:fm1}, right, 
include 
{\em (i)} wire-wrapped fuel-rod bundles, 
{\em (ii)} an inertial-confinement fusion (ICF) miniapp, and 
{\em (iii)} heat-transfer enhancement through optimized wire-coil-inserts.
  These components exemplify systems central to DOE’s advanced energy
programs and are designed to permit ready exploration, development, and
validation of agentic workflows in this space, where
    energy transfer is characterized by complex turbulent flow, strong
geometric sensitivity, and stringent performance and safety constraints.
  They also constitute a rich data domain, combining decades of high-fidelity
simulations using Nek5000/RS (one of the few CFD codes accepted in Nuclear
Regulatory Commission (NRC)
licensing workflows and a former Gordon Bell Prize winner
\cite{tufo99a,gb23}) and MFEM (the 2025 Gordon Bell Winner \cite{gb25mfem}) 
with extensive experimental datasets from both open literature and proprietary
industry sources.  These two teams have worked closely together in the last
decade to deliver major application and HPC advancements to the DOE under the
Exascale Computing Project \cite{ceed}.  

%  The central objective of this {\bf Phase I effort} is to {\em unify the
%  available heterogeneous datasets into a common (mesh-based) geometry-aware
%  representation and to demonstrate rapid, accurate, surrogate solution}
%  generation that will be applied to design and optimization of thermal-fluid
%  systems.

The project team includes leading experts (Merzari and Shaver) in fission and
fusion thermal-hydraulics who will ensure that training and quantities of
interest (QOIs) are directly relevant to fusion/fission design work at national
laboratories and in industry, where we are partnering with several firms,
including {\bf Kairos}, {\bf TerraPower}, and {\bf Westinghouse} 
             (Appendix \ref{collaboration_letters}).
In these systems, fluid motion is the {\em principal mechanism} of
energy transfer from core processes (fission/fusion reactions) to heat exchangers, where
recipient working fluids ultimately drive turbines or other power-generation
devices.
%
Combustion systems and inertially confined fusion (ICF) similarly depend on
controlled fluid motion to initiate and sustain energy-producing reactions.
Computational costs are significant because of complex geometries, broad ranges
of timescales (e.g., convective flow versus thermal diffusion in solids), and
required accuracy for safety.
  Current {\bf system- and device-level optimization—particularly when involving
geometric or topology modifications—can require tens to hundreds of
high-fidelity simulations}, each taking days to months on high-performance
computing (HPC) platforms, with {\bf no shared information across the
simulation space.}   

  This Phase I effort will combine graph neural networks (GNNs) with
autoencoders to build a {\bf common latent space} compressing
geometric, fluid, and thermal information {\bf across all project data sets.}
This shared information space will permit significantly shortened training
cycles for {\bf new device configurations} (e.g., fuel rods and heat
exchangers).  Surrogate prediction quality---with emphasis on
design-driven QOIs---will be assessed directly through comparisons with
available experimental
data and current simulation capabilities.  The ultimate goal, to be realized in
Phase II, is to radically transform the design cycle for energy-system
components by providing reliable, high-accuracy AI-driven surrogate models
tightly coupled with state-of-the-art design optimization methods.  Runtime
surrogate uncertainty quantification (UQ) will be a key component of Phase II
to identify necessary retraining (where UQ is high) at low cost.
% so that necessary retraining (where UQ is high) may be identified at low cost.

\section{Project Objectives}

\input objective

\input proposed_methods

\input milestone

\input data

\input metric
